\chapter[Ultraviolet completion]{Ultraviolet Operators, Ghosts, and a Cubic Completion}
\label{ch:ultraviolet}

\paragraph{Learning outcomes.}
After this chapter the reader should be able to:
\begin{enumerate}
  \item read the \(a_2\) density as a list of four-derivative
        operators and assign each a role (correction versus
        propagator);
  \item state Ostrogradsky's theorem in the form used for
        gravitational effective field theory;
  \item explain why quadratic gravity of Stelle type is not an
        acceptable infrared limit of NAMCG;
  \item write the Einsteinian cubic invariant and the property that
        makes it safer than a generic \(R^3\) term;
  \item enter Chapter~10 knowing that the ultraviolet problem of
        the metric effective action is under a named control,
        not solved in the sense of a non-perturbative continuum
        limit.
\end{enumerate}

Chapter~8 produced
\begin{equation}
  \label{eq:W-again-9}
  W[g]
  =\int\mathrm{d}^4x\,\sqrt{-g}
   \Bigl(
     -\rho_\Lambda
     +\frac{R}{16\pi\Geff}
     +\alpha\,a_2[g]
     +\cdots
   \Bigr).
\end{equation}
The first two terms are the cosmological debt and the Einstein--Hilbert
term. This chapter treats the third.

\section{What \(a_2\) contains}
\label{sec:ch9-a2}

On a Laplace-type bundle the second Seeley--DeWitt coefficient is a
linear combination of
\begin{equation}
  \label{eq:a2-ops}
  R_{\mu\nu\rho\sigma}R^{\mu\nu\rho\sigma},
  \quad
  R_{\mu\nu}R^{\mu\nu},
  \quad
  R^2,
  \quad
  E^2,
  \quad
  RE,
  \quad
  D^2 R,\quad
  D^2 E,
\end{equation}
together with bundle-curvature terms \(F_{\mu\nu}F^{\mu\nu}\). After
the endomorphism \(E\) is evaluated on the condensate and total
derivatives are discarded, the metric operators that remain in four
dimensions may be written
\begin{equation}
  \label{eq:W4}
  W_4[g]
  =\int\mathrm{d}^4x\,\sqrt{-g}
   \bigl(
     c_W C_{\mu\nu\rho\sigma}C^{\mu\nu\rho\sigma}
     +c_R R^2
     +c_{\mathrm{GB}}\,E_{\mathrm{GB}}
   \bigr),
\end{equation}
where \(C_{\mu\nu\rho\sigma}\) is the Weyl tensor and
\begin{equation}
  \label{eq:GB}
  E_{\mathrm{GB}}
  =R_{\mu\nu\rho\sigma}R^{\mu\nu\rho\sigma}
   -4R_{\mu\nu}R^{\mu\nu}
   +R^2
\end{equation}
is the Gauss--Bonnet density, topological in four dimensions at the
level of the bulk equations of motion.

The coefficients \(c_W,c_R,c_{\mathrm{GB}}\) are logarithms of
\(\LUV/\mu\) times fibre traces, as in~\eqref{eq:W-expand}. Their
precise rationals are scheme-dependent in the same sense as \(A,B,C\)
in~\eqref{eq:Geff-structure}. What is not scheme-dependent is the
existence of the Weyl-squared and \(R^2\) operators at this order.

\section{Ostrogradsky's theorem, used as a veto}
\label{sec:ch9-ostro}

\begin{theorem}[Ostrogradsky, gravitational form]
  \label{thm:ostro}
  Let a field theory possess a Lagrangian that depends on second
  (or higher) time derivatives of a dynamical variable, not
  removable by adding a total derivative or by imposing a
  degeneracy condition that reduces the order. Then the Hamiltonian
  of that variable is unbounded below. In a quadratic gravitational
  action about Minkowski spacetime, the unbounded direction is a
  ghost pole in the propagator of \(h_{\mu\nu}\).
\end{theorem}

The theorem is the reason one may not promote~\eqref{eq:W4} to a
fundamental action for a propagating metric and call the result a
completion of NAMCG.

\paragraph{Stelle quadratic gravity.}
The action
\begin{equation}
  \label{eq:stelle}
  S_{\mathrm{Stelle}}
  =\int\mathrm{d}^4x\,\sqrt{-g}
   \Bigl(
     \frac{R}{16\pi G}
     -\frac{\alpha}{2}C_{\mu\nu\rho\sigma}C^{\mu\nu\rho\sigma}
     +\frac{\beta}{2}R^2
   \Bigr)
\end{equation}
propagates, in addition to the massless graviton, a massive spin-2
ghost and a massive scalar. The ghost mass is of order
\(1/\sqrt{\alpha G}\). If \(\alpha\) is a one-loop logarithm, that
mass is near \(\LUV\), which is the Planck scale of
Section~\ref{sec:ch8-size}. A Planck-scale ghost is an ultraviolet
sickness of the metric effective theory. It is not an infrared
sickness provided those poles are never treated as asymptotic
states of the condensate.

\paragraph{Effective-field-theory reading.}
The same operators are harmless if they are used only as vertices
in a derivative expansion of matrix elements with external momenta
\(p\ll\LUV\), and never as additional poles. That is the reading
adopted here:

\begin{assumption}[Four-derivative terms are EFT vertices]
  \label{ass:EFT-a2}
  The operators in~\eqref{eq:W4} correct graviton scattering and
  static potentials at order \((p/\LUV)^2\). They are not to be
  diagonalized into extra propagating modes of \(h_{\mu\nu}\).
\end{assumption}

Assumption~\ref{ass:EFT-a2} is available only because the metric is
composite. The microscopic path integral is over \(\Phi\), not over
an independent \(g_{\mu\nu}\) with action~\eqref{eq:stelle}. The
poles of Stelle gravity are an artefact of treating \(W[g]\) as a
fundamental metric action beyond its radius of convergence.

\section{What ultraviolet stability does mean}
\label{sec:ch9-stability}

Presentation~A is a matrix model. Its ultraviolet definition is the
regularized trace over \(\mathcal{H}\) with cutoff~\eqref{eq:LUV-def}.
Questions that make sense at that level are:

\begin{itemize}
  \item boundedness of \(S_{\mathrm{matrix}}\) from below on Hermitian
        \(\hat X^\mu\);
  \item existence of the condensed saddle
        \(\langle\Phi^I{}_A\rangle=v\Lambda^A{}_I\);
  \item absence of tachyons in the blocks that Theorem~\ref{thm:mT}
        and Assumption~\ref{ass:gap-psi} did not already constrain.
\end{itemize}

Questions that do \emph{not} make sense until a continuum limit in
\(\LUV\) is constructed are:

\begin{itemize}
  \item renormalizability of \(W[g]\) as a fundamental quantum field
        theory of \(g_{\mu\nu}\);
  \item asymptotic safety of that metric theory;
  \item unitarity of S-matrix elements among Stelle poles.
\end{itemize}

NAMCG, as written in this book, claims the first list and refuses
the second. The metric effective action is an infrared object with
a cutoff. Ultraviolet stability means stability of the matrix
saddle, not renormalizability of quadratic gravity.

A sufficient condition for the matrix action~\eqref{eq:Smatrix} to
be bounded below on Hermitian matrices is that the quartic
commutator term dominates with the sign written there,
\begin{equation}
  \label{eq:YM-bounded}
  -\mathrm{Tr}\bigl([\hat X^\mu,\hat X^\nu][\hat X_\mu,\hat X_\nu]\bigr)
  \ge 0
\end{equation}
in the Euclideanized internal-time convention of IKKT-type models,
together with a potential \(V\) that is bounded below, as
\eqref{eq:V} is. Lorentzian signature complicates the bound because
\(\hat X^0\) enters with the opposite sign in some reductions. That
complication is inherited from the matrix-model literature and is
not solved here. It is recorded in Chapter~17.

\section{Cubic operators and Einsteinian cubic gravity}
\label{sec:ch9-ECG}

Singularities of the Einstein equation are not cured by the
four-derivative vertices of Assumption~\ref{ass:EFT-a2}. Those
vertices are small when curvature is small, and they are
untrustworthy when curvature reaches \(\LUV^2\). A different
higher-curvature object is needed if compact-object cores are to
be regularized inside the derivative expansion.

Einsteinian cubic gravity (ECG) is the four-dimensional cubic
curvature invariant whose linearized spectrum about a maximally
symmetric background coincides with that of Einstein gravity: no
extra spin-2 ghost, no extra scalar at that order. A convenient
writing is
\begin{equation}
  \label{eq:ECG}
  \mathcal{P}
  =12 R_{\mu}{}^{\rho}{}_{\nu}{}^{\sigma}
       R_{\rho}{}^{\tau}{}_{\sigma}{}^{\eta}
       R_{\tau}{}^{\mu}{}_{\eta}{}^{\nu}
   +\cdots,
\end{equation}
the omitted terms being the unique combination of
\(R^3\), \(RR_{\mu\nu}R^{\mu\nu}\), and
\(R_{\mu\nu\rho\sigma}R^{\mu\nu\rho\sigma}R\) that implements the
spectral condition. The action used later in the book is
\begin{equation}
  \label{eq:S-ECG}
  S_{\mathrm{ECG}}[g]
  =\int\mathrm{d}^4x\,\sqrt{-g}
   \Bigl(
     \frac{R-2\Lambda_{\mathrm{eff}}}{16\pi\Geff}
     +\beta_{\mathrm{ECG}}\,\mathcal{P}
   \Bigr).
\end{equation}

Two properties matter.

\begin{enumerate}
  \item For the first-integral equation displayed in
        Section~\ref{sec:status-J}, the static spherical equation is
        a nonlinear second-order ODE for \(f\), not an algebraic
        prescription; that this equation follows from~\eqref{eq:S-ECG}
        remains to be derived. The
        regular-core profile
        \begin{equation}
          \label{eq:fr-ECG}
          f(r)
          =1-\frac{2M r^2}{r^3+r_0^3}
           +\frac{\Lambda_{\mathrm{eff}}}{3}r^2
        \end{equation}
        used in Chapter~13 is a phenomenological regular-core model
        with \(r_0^3\propto\beta_{\mathrm{ECG}}M\). For
        \(\Lambda_{\mathrm{eff}}=0\) and \(Mr_0\ne0\) it is not an
        exact vacuum solution of the first-integral equation of
        Section~\ref{sec:status-J}
        (conditional on that equation following from~\eqref{eq:S-ECG}).
  \item The spectral condition is a statement about linearization
        on maximally symmetric backgrounds. It is not a proof that
        every time-dependent perturbation of~\eqref{eq:fr-ECG} is
        ghost-free. The book uses ECG as a controlled model of
        core regularity, not as a claim of non-perturbative
        unitarity.
\end{enumerate}

\begin{assumption}[Cubic completion]
  \label{ass:ECG}
  Through the curvature scales relevant to compact-object cores,
  the metric effective action may be truncated at
  \eqref{eq:S-ECG}, with \(\beta_{\mathrm{ECG}}\) set by the
  matching of \(a_2\) and by the matrix scale \(r_0\sim\LUV^{-1}\)
  up to an order-one factor.
\end{assumption}

Assumption~\ref{ass:ECG} is stronger than Assumption~\ref{ass:EFT-a2}.
It promotes one cubic combination from a vertex to a term that is
solved non-perturbatively in the static spherical sector. The
promotion is justified only in that sector, and only as a model.

\section{Torsion and higher connections}
\label{sec:ch9-torsion}

A chosen non-degenerate frame (Repair~II,~\eqref{eq:vev-is-frame}, or an
assumed coframe) need not have
vanishing torsion. A torsion tensor
\begin{equation}
  \label{eq:torsion}
  T^A{}_{\mu\nu}
  =\partial_\mu e^A_\nu-\partial_\nu e^A_\mu
   +\omega_\mu{}^A{}_B e^B_\nu
   -\omega_\nu{}^A{}_B e^B_\mu
\end{equation}
is an additional infrared field if \(\omega\) is independent of
\(e\). In NAMCG the connection is not independent: it is the
connection that eats the Goldstone \(\omega_{AB}\) of Chapter~5.
After that eating, torsion is either zero or massive of order
\(v\). Massive torsion is an ultraviolet correction of the same
type as the \(a_2\) vertices. It is not retained as a propagating
infrared field.

\section{A checklist before Lorentz locking}
\label{sec:ch9-check}

The metric effective theory, as it now stands, is:

\begin{center}
\begin{tabular}{@{}ll@{}}
\toprule
Ingredient & Status \\
\midrule
massless TT tensor & Thm.~\ref{thm:mT} \\
prefactor \(\Geff\) & Thm.~\ref{thm:Geff} \\
quartic vacuum energy & isolated, not tamed \\
\(C^2\) and \(R^2\) & EFT vertices, Assp.~\ref{ass:EFT-a2} \\
cubic core regularity & ECG model, Assp.~\ref{ass:ECG} \\
torsion & massive or zero after eating \\
local Lorentz invariance as a gauge symmetry
  of \(\Phi\) & not yet constructed \\
\bottomrule
\end{tabular}
\end{center}

The missing gauge symmetry is the last geometric ingredient. Without
it the identification of \(\omega_{AB}\) as a Goldstone, used since
Chapter~5, is incomplete.

\section{Exercises}
\label{sec:ch9-exercises}

\begin{enumerate}
  \item \textbf{Gauss--Bonnet.}
    Show that \(E_{\mathrm{GB}}\) in four dimensions is a total
    derivative plus a topological density. Why, then, can
    \(c_{\mathrm{GB}}\) still matter on manifolds with boundary?
  \item \textbf{Stelle poles.}
    Linearize~\eqref{eq:stelle} about Minkowski spacetime in
    harmonic gauge. Locate the massive spin-2 pole and determine
    the sign of its residue relative to the massless graviton.
  \item \textbf{EFT versus poles.}
    A tree-level graviton-graviton amplitude receives a correction
    \(\sim c_W s^2/\LUV^2\) from~\eqref{eq:W4}. At what
    centre-of-mass energy does this correction become unreliable
    as an EFT vertex, and why is that energy the same scale at
    which Stelle poles would have appeared?
  \item \textbf{Bounded matrix action.}
    For two Hermitian matrices \(X,Y\) show that
    \(-\mathrm{Tr}([X,Y]^2)\ge 0\). Identify the step that fails,
    or requires a Euclidean rotation, when one matrix is
    \(X^0\) with Lorentzian signature.
  \item \textbf{Spherical reduction.}
    Take \(ds^2=-f(r)dt^2+f(r)^{-1}dr^2+r^2 d\Omega^2\) and list
    which curvature scalars in a generic \(R^3\) combination
    survive. The displayed first integral of
    Section~\ref{sec:status-J} is a second-order equation in \(f\),
    not a fourth-order one. Verify the leading \(1/r\) terms of that
    equation as an asymptotic consistency check. This does not derive
    the equation from~\(\mathcal{P}\); that derivation is the open
    task of Problem~J.
  \item \textbf{Torsion mass.}
    Give a dimensional estimate of the torsion mass if the only
    scale available is \(v\). Compare it with \(m_\psi\) from
    Assumption~\ref{ass:gap-psi}.
\end{enumerate}

\section{Bridge to Chapter~10}
\label{sec:ch9-bridge}

The infrared metric theory is Einstein gravity plus controlled
higher-curvature corrections and an untamed \(\rho_\Lambda\). What
it still lacks is an explicit demonstration that local
\(\mathrm{SO}(3,1)\) is a gauge symmetry of the condensate and that
the six Goldstones \(\omega_{AB}\) are eaten by the connection,
leaving no extra massless scalars. Chapter~10 constructs that
locking. After it, the counting table of Section~\ref{sec:ch5-count}
is complete through the line ``eat \(\omega_{AB}\),'' and the
frame, taken from Repair~II,~\eqref{eq:vev-is-frame} (or from an
assumed non-degenerate coframe), transforms as a vierbein under the
diagonal local Lorentz group, under the assumptions of
Theorem~\ref{thm:locking}. Definition~\ref{def:vierbein} is not
used to construct it (Problem~C).